\scanpage{7}
% source lines 1-2, continuation of scan 6
de diviseurs absolus \inserted{(de \deleted{norme} valeur absolue $1$)} est compact. Or on a
\[
{\unsure{D|\mathscr{D}}}_{\mathrm{abs}}(K)\simeq D(K)\times D_\infty(K)
\]

% source lines 3-5
\deleted{où $D(K)$ est le \unsure{gr.}}
\[
\left\{\begin{aligned}
&D(K)\quad\text{groupe des diviseurs usuels de }K\\
&D_\infty(K)=\coprod_{\mathfrak p\text{ à l'infini}}W_{\mathfrak p}\quad\text{groupe des diviseurs à l'infini}
\end{aligned}\right.
\]

% source lines 6-8
On \deleted{$\unsure{K}$} écrira
\[
0\longrightarrow D_{\infty,0}(K)\longrightarrow D_{\mathrm{abs},0}(K)\longrightarrow D(K)\longrightarrow0
\]
et on considère dans $D_{\mathrm{abs},0}(K)$ l'image de $K^*$, soit $P\simeq K^*/\unsure{\mu|U}(K)$ \deleted{\unsure{des}} diviseurs absolus principaux.

On a alors
\[
0\longrightarrow D_{\infty,0}(K)/E_0(K)\longrightarrow{\unsure{\mathscr{P}|\mathfrak{P}|\mathcal{P}}}_{\mathrm{abs},0}(K)\longrightarrow\mathscr{P}(K)\longrightarrow0.
\]

% source lines 10-14
où on pose
\[
\begin{aligned}
\bigl(\quad&\mathscr{P}(K)=D(K)/\operatorname{Im}K^*=D_{\mathrm{abs}}(K)/P+D_\infty(K)\\
\bigl(\quad&E_0(K)\simeq P\cap D_\infty^{\deleted{\unsure{0}}}(K)\simeq(P+D_\infty(K))/P\simeq\deleted{\unsure{E(K)}}E(K)/\mu(K),\ \text{\unsure{où}}.
\end{aligned}
\]
$\bigl(\ E(K)\simeq$ unités de $K$ i.e. \unsure{élmts|éléments|élts} de $K^*$ dont l'image dans $D_{\mathrm{abs}}(K)$ est dans $D_\infty(K)$.

% source lines 15-16
\deleted{\unsure{D'autre part, le}} Ceci dit, le th. de \deleted{Minkowski} \inserted{Dirichlet} \unsure{équivaut} à \unsure{l'ens|l'un} des deux corollaires

% source lines 17-20
\uline{Cor. 1} $\mathscr{P}(K)$ est fini

\uline{Cor. 2} $E_0(K)$ est de rang $r_1+r_2-1$

(i.e. la dimension de $D_{\infty,0}(K)$). D'où

\uline{Cor. 3} $\mathscr{P}_{\mathrm{abs},0}(K)$ est une extension de $\mathscr{P}(K)$, \unsure{qui est} un \unsure{grp|gpe} fini, par un tore de dim. $r_1+r_2-1$.
% archive sheet number: 6
